\documentclass[a4paper,12pt]{article}

% -------------------------------------------------
% Кодировка и русский язык
% -------------------------------------------------
\usepackage[T2A]{fontenc}
\usepackage[utf8]{inputenc}
\usepackage[russian]{babel}

% -------------------------------------------------
% Математика и шрифт
% -------------------------------------------------
\usepackage{amsmath,amssymb,amsthm,mathtools}
\usepackage{helvet}
\renewcommand{\familydefault}{\sfdefault}
\usepackage{icomma}

% -------------------------------------------------
% Вёрстка
% -------------------------------------------------
\usepackage[
  left=20mm,
  right=20mm,
  top=19mm,
  bottom=20mm
]{geometry}
\usepackage{microtype}
\usepackage{setspace}
\onehalfspacing
\usepackage{parskip}
\setlength{\parindent}{0pt}

% -------------------------------------------------
% Таблицы и списки
% -------------------------------------------------
\usepackage{tabularx,booktabs}
\usepackage{enumitem}
\setlist[itemize]{
  leftmargin=6mm,
  itemsep=2pt,
  topsep=3pt
}
\setlist[enumerate]{
  leftmargin=8mm,
  itemsep=6pt,
  topsep=4pt
}
\renewcommand{\arraystretch}{1.3}

% -------------------------------------------------
% Графика и оформление
% -------------------------------------------------
\usepackage{tikz}
\usetikzlibrary{arrows.meta,positioning}
\usepackage{xcolor}
\usepackage{titlesec}
\usepackage{fancyhdr}
\usepackage{hyperref}

\definecolor{accent}{HTML}{008080}
\definecolor{darkaccent}{HTML}{004D4D}
\definecolor{textgray}{HTML}{333333}
\definecolor{softgray}{HTML}{F1F4F4}
\definecolor{linegray}{HTML}{AAB8B8}
\definecolor{warning}{HTML}{8A4B3A}

\color{textgray}

\hypersetup{
  colorlinks=true,
  linkcolor=darkaccent,
  urlcolor=darkaccent,
  pdftitle={Вклады: сложные проценты и промежуточные операции},
  pdfauthor={Лёвин Артём Александрович}
}

\titleformat{\section}
  {\Large\bfseries\color{darkaccent}}
  {\thesection.}
  {0.5em}
  {}
  [\vspace{1mm}\color{accent}\titlerule]

\titleformat{\subsection}
  {\large\bfseries\color{accent}}
  {\thesubsection.}
  {0.5em}
  {}

\titlespacing*{\section}{0pt}{13pt}{7pt}
\titlespacing*{\subsection}{0pt}{10pt}{4pt}

\pagestyle{fancy}
\fancyhf{}
\fancyhead[L]{\small\color{textgray}Вклады и сложные проценты}
\fancyhead[R]{\small\color{textgray}Тимофей Васильченко}
\fancyfoot[C]{\small\thepage}
\renewcommand{\headrulewidth}{0.4pt}
\renewcommand{\headrule}{
  \hbox to\headwidth{\color{linegray}\leaders\hrule height \headrulewidth\hfill}
}

\newcommand{\rub}{\,\text{руб.}}
\newcommand{\RUB}{\text{руб.}}
\newcommand{\important}[1]{\textbf{\color{darkaccent}#1}}
\newcommand{\warningtext}[1]{\textbf{\color{warning}#1}}
\newcommand{\formula}[1]{%
  \[
    \boxed{\displaystyle #1}
  \]
}

\begin{document}

% =================================================
% Титульный лист
% =================================================
\begin{titlepage}
\thispagestyle{empty}
\centering

\vspace*{28mm}

{\Large\color{accent}\textbf{ЧЕК-ЛИСТ}\par}

\vspace{8mm}

{\huge\bfseries\color{darkaccent}
Вклады: сложные проценты\\[2mm]
и промежуточные операции
\par}

\vspace{13mm}

{\large
Математическое моделирование финансовых задач
\par}

\vspace{24mm}

{\large
Лёвин Артём Александрович\\[2mm]
эксклюзивно для Тимофея Васильченко
\par}

\vspace{12mm}

{\large \today\par}

\vfill

\begin{tikzpicture}[x=1cm,y=1cm]
  \draw[
    color=accent!55,
    line width=1.1pt
  ]
  (0,0)
  .. controls (2.2,0.75) and (4.1,-0.55) .. (6.8,0.10)
  .. controls (9.2,0.70) and (11.6,-0.35) .. (14,0.25);
\end{tikzpicture}

\vspace{7mm}

\end{titlepage}

% =================================================
% Основной чек-лист
% =================================================
\section{Главная идея банковского вклада}

Вкладчик размещает некоторую сумму в банке. По окончании каждого расчётного
периода банк начисляет процент на сумму, находившуюся на счёте к моменту
начисления.

В рамках данного чек-листа предполагается, что:

\begin{itemize}
  \item расчётный период равен одному году;
  \item проценты начисляются один раз в конце каждого года;
  \item дополнительные взносы и снятия выполняются \important{после начисления процентов};
  \item процентная ставка в течение всего срока не изменяется.
\end{itemize}

Если на счёте находилось \(100\RUB\), а ставка равна \(20\%\), то за год банк
начислит
\[
100\cdot\frac{20}{100}=20\RUB.
\]
После начисления процентов на счёте будет
\[
100+20=120\RUB.
\]

\subsection{Последовательность операций}

\begin{center}
\begin{tikzpicture}[
  node distance=6mm,
  box/.style={
    draw=accent,
    rounded corners=2pt,
    minimum height=11mm,
    text width=31mm,
    align=center,
    line width=0.8pt,
    fill=softgray
  },
  arrow/.style={
    -{Latex[length=2.4mm]},
    color=darkaccent,
    line width=0.8pt
  }
]
\node[box] (start) {Сумма в начале года};
\node[box, right=of start] (percent) {Начисление процентов};
\node[box, right=of percent] (operation) {Взнос или снятие};
\node[box, right=of operation] (finish) {Сумма в конце года};

\draw[arrow] (start) -- (percent);
\draw[arrow] (percent) -- (operation);
\draw[arrow] (operation) -- (finish);
\end{tikzpicture}
\end{center}

\section{Основные обозначения}

Пусть:

\begin{itemize}
  \item \(S_0\) --- первоначальная сумма вклада;
  \item \(S_k\) --- сумма на счёте после завершения \(k\)-го года;
  \item \(P\%\) --- годовая процентная ставка;
  \item \(p\) --- ставка, записанная в виде десятичной дроби;
  \item \(r\) --- коэффициент увеличения суммы за один год;
  \item \(x_k\) --- дополнительная операция после начисления процентов
  в \(k\)-м году.
\end{itemize}

Процентную ставку переводят в десятичную дробь:

\formula{p=\frac{P}{100}}

Коэффициент ежегодного увеличения:

\formula{r=1+p=1+\frac{P}{100}}

Например, при ставке \(9\%\):
\[
p=\frac{9}{100}=0{,}09,
\qquad
r=1+0{,}09=1{,}09.
\]

\important{Не путайте величины:}
\[
P=9,\qquad p=0{,}09,\qquad r=1{,}09.
\]

\section{Базовая таблица финансовой задачи}

Для каждого года удобно заполнять четыре столбца:

\begin{table}[ht]
\centering
\caption{Структура математической модели вклада}
\begin{tabularx}{\linewidth}{@{}cXXXX@{}}
\toprule
Год &
Сумма в начале года &
Начисленные проценты &
Дополнительная операция &
Сумма в конце года \\
\midrule
\(k\) &
\(S_{k-1}\) &
\(S_{k-1}p\) &
\(x_k\) &
\(S_{k-1}+S_{k-1}p+x_k\) \\
\bottomrule
\end{tabularx}
\end{table}

Так как
\[
S_{k-1}+S_{k-1}p
=
S_{k-1}(1+p)
=
S_{k-1}r,
\]
получаем основную рекуррентную формулу:

\formula{S_k=S_{k-1}r+x_k}

Знаки дополнительных операций:

\[
x_k>0
\quad\Longleftrightarrow\quad
\text{деньги внесли на счёт},
\]
\[
x_k<0
\quad\Longleftrightarrow\quad
\text{деньги сняли со счёта}.
\]

Если после начисления процентов сняли \(3000\RUB\), то
\[
x_k=-3000.
\]

\section{Вклад без дополнительных операций}

Если деньги не вносят и не снимают, то
\[
x_1=x_2=\ldots=x_n=0.
\]

После первого года:
\[
S_1=S_0r.
\]

После второго года:
\[
S_2=S_1r=S_0r^2.
\]

После третьего года:
\[
S_3=S_2r=S_0r^3.
\]

Следовательно, после \(n\) лет:

\formula{S_n=S_0r^n
=S_0\left(1+\frac{P}{100}\right)^n}

Это \important{формула сложных процентов}.

Проценты называются сложными, потому что каждый следующий процент начисляется
не только на первоначальную сумму, но и на проценты, начисленные ранее.

\subsection{Пример}

На вклад положили \(100\RUB\) под \(20\%\) годовых на три года.

Здесь
\[
r=1+\frac{20}{100}=1{,}2.
\]

\begin{table}[ht]
\centering
\caption{Изменение суммы вклада}
\begin{tabularx}{0.9\linewidth}{@{}cXXXX@{}}
\toprule
Год &
Начальная сумма &
Проценты &
Операция &
Конечная сумма \\
\midrule
1 & \(100\) & \(20\) & \(0\) & \(120\) \\
2 & \(120\) & \(24\) & \(0\) & \(144\) \\
3 & \(144\) & \(28{,}8\) & \(0\) & \(172{,}8\) \\
\bottomrule
\end{tabularx}
\end{table}

Проверка по формуле:
\[
S_3=100\cdot1{,}2^3=172{,}8\RUB.
\]

\section{Дополнительные взносы и снятия}

Если после каждого начисления процентов выполняется дополнительная операция,
необходимо последовательно применять формулу
\[
S_k=S_{k-1}r+x_k.
\]

\subsection{Пример с разными операциями}

На счёт положили \(100\RUB\) под \(20\%\) годовых.

\begin{itemize}
  \item после первого года внесли \(10\RUB\);
  \item после второго года внесли \(10\RUB\);
  \item после третьего года сняли \(20\RUB\).
\end{itemize}

Тогда:

\[
S_1=100\cdot1{,}2+10=130;
\]
\[
S_2=130\cdot1{,}2+10=166;
\]
\[
S_3=166\cdot1{,}2-20=179{,}2.
\]

\begin{table}[ht]
\centering
\caption{Вклад с промежуточными операциями}
\begin{tabularx}{0.9\linewidth}{@{}cXXXX@{}}
\toprule
Год &
Начальная сумма &
Проценты &
Операция &
Конечная сумма \\
\midrule
1 & \(100\) & \(20\) & \(+10\) & \(130\) \\
2 & \(130\) & \(26\) & \(+10\) & \(166\) \\
3 & \(166\) & \(33{,}2\) & \(-20\) & \(179{,}2\) \\
\bottomrule
\end{tabularx}
\end{table}

\section{Общая формула с промежуточными операциями}

Пусть после первого года выполняется операция \(x_1\), после второго ---
\(x_2\), и так далее.

Тогда:
\[
S_1=S_0r+x_1,
\]
\[
S_2=(S_0r+x_1)r+x_2
=S_0r^2+x_1r+x_2,
\]
\[
S_3=(S_0r^2+x_1r+x_2)r+x_3
=S_0r^3+x_1r^2+x_2r+x_3.
\]

После \(n\) лет:

\formula{
S_n
=
S_0r^n
+x_1r^{n-1}
+x_2r^{n-2}
+\ldots
+x_{n-1}r
+x_n
}

Степень коэффициента \(r\) показывает, сколько лет соответствующая сумма
находилась на счёте после внесения.

\begin{itemize}
  \item Первоначальная сумма \(S_0\) работала все \(n\) лет, поэтому умножается
  на \(r^n\).
  \item Сумма \(x_1\), внесённая после первого года, работала ещё \(n-1\) лет.
  \item Последняя сумма \(x_n\) вносится после последнего начисления и уже
  не приносит процентов.
\end{itemize}

\section{Одинаковый взнос после каждого года}

Если после каждого года вносится одна и та же сумма \(x\), то
\[
x_1=x_2=\ldots=x_n=x.
\]

Получаем:
\[
S_n
=
S_0r^n
+xr^{n-1}
+xr^{n-2}
+\ldots
+xr+x.
\]

Или
\[
S_n
=
S_0r^n+x\left(r^{n-1}+r^{n-2}+\ldots+r+1\right).
\]

Сумма в скобках является суммой геометрической прогрессии. При \(r\ne1\):

\formula{
S_n=S_0r^n+x\frac{r^n-1}{r-1}
}

Так как \(r-1=p\), формулу также можно записать в виде
\[
S_n=S_0r^n+x\frac{r^n-1}{p}.
\]

\warningtext{Ограничение:} делить на \(r-1\) можно только при \(r\ne1\),
то есть при ненулевой процентной ставке.

Если \(P=0\), то \(r=1\), и сумма вычисляется отдельно:
\[
S_n=S_0+nx.
\]

\section{Через сколько лет вклад увеличится в несколько раз}

Пусть требуется определить, через сколько лет вклад станет в \(k\) раз больше
первоначального.

Без дополнительных операций:
\[
S_n=S_0r^n.
\]

По условию:
\[
S_n=kS_0.
\]

Следовательно,
\[
S_0r^n=kS_0.
\]

При \(S_0>0\) можно разделить обе части на \(S_0\):
\[
r^n=k.
\]

Логарифмируем обе части:
\[
\lg r^n=\lg k.
\]

Используем свойство
\[
\lg r^n=n\lg r.
\]

Отсюда:

\formula{n=\frac{\lg k}{\lg r}}

Условия применимости формулы:
\[
r>0,\qquad r\ne1,\qquad k>0.
\]

Для обычной задачи о росте вклада:
\[
P>0,\qquad r>1,\qquad k>1.
\]

\subsection{Пример: удвоение вклада}

Ставка равна \(9\%\). Требуется определить, через сколько полных лет вклад
увеличится не менее чем в два раза.

\[
r=1+\frac{9}{100}=1{,}09.
\]

Решаем уравнение:
\[
1{,}09^n=2.
\]

\[
n=\frac{\lg2}{\lg1{,}09}\approx8{,}04.
\]

Через \(8\) полных лет вклад ещё не достигнет требуемой величины. Поэтому
необходимо округлить результат \important{в большую сторону}:

\[
n_{\min}=\left\lceil8{,}04\right\rceil=9.
\]

Ответ: вклад увеличится не менее чем в два раза через \(9\) полных лет.

\section{Смысловое округление}

В финансовых задачах результат округляется не только по обычным
арифметическим правилам, но и с учётом смысла вопроса.

Если вычисления дали
\[
n=8{,}13,
\]
это означает, что через \(8\) полных лет нужный результат ещё не достигнут.

Если требуется найти минимальное целое количество лет, за которое сумма станет
не меньше заданной, используют округление вверх:
\[
n_{\min}=\lceil n\rceil.
\]

\important{Контрольный вопрос:} выполнено ли требуемое условие после найденного
целого числа периодов?

\section{Сравнение запланированной и фактической суммы}

Часто требуется сравнить:

\begin{itemize}
  \item сумму без промежуточных операций;
  \item сумму с учётом снятий и дополнительных взносов.
\end{itemize}

Разность удобно обозначать буквой \(\Delta\):
\[
\Delta
=
S_{\text{без операций}}
-
S_{\text{с операциями}}.
\]

Если \(\Delta>0\), фактическая сумма меньше запланированной на
\(\Delta\RUB\).

\subsection{Типовой пример}

Вклад открыт под \(10\%\) годовых на три года.

\begin{itemize}
  \item после первого года сняли \(2000\RUB\);
  \item после второго года внесли \(2000\RUB\);
  \item после третьего года операций не было.
\end{itemize}

Пусть первоначальная сумма равна \(S\), а
\[
r=1{,}1.
\]

Без операций:
\[
S_{\text{без операций}}=Sr^3.
\]

С операциями:
\[
S_{\text{с операциями}}
=
\bigl((Sr-2000)r+2000\bigr)r.
\]

Раскрываем скобки:
\[
S_{\text{с операциями}}
=
Sr^3-2000r^2+2000r.
\]

Находим разность:
\begin{align*}
\Delta
&=
Sr^3-\left(Sr^3-2000r^2+2000r\right)\\
&=
2000r^2-2000r\\
&=
2000r(r-1).
\end{align*}

Подставляем \(r=1{,}1\):
\[
\Delta
=
2000\cdot1{,}1\cdot0{,}1
=
220\RUB.
\]

Первоначальная сумма \(S\) сократилась, поэтому для решения задачи её значение
не требовалось.

\section{Универсальный алгоритм решения}

\begin{enumerate}
  \item Определите первоначальную сумму, ставку, количество периодов и моменты
  выполнения дополнительных операций.

  \item Переведите процентную ставку в десятичную дробь:
  \[
  p=\frac{P}{100}.
  \]

  \item Введите коэффициент:
  \[
  r=1+p.
  \]

  \item Зафиксируйте порядок событий внутри каждого года:
  \[
  \text{начальная сумма}
  \longrightarrow
  \text{проценты}
  \longrightarrow
  \text{операция}.
  \]

  \item Составьте таблицу по годам.

  \item Используйте переход:
  \[
  S_k=S_{k-1}r+x_k.
  \]

  \item Снятие записывайте со знаком минус, внесение --- со знаком плюс.

  \item Выносите повторяющиеся выражения за скобки и только затем раскрывайте
  скобки, если это требуется.

  \item Составьте уравнение или разность, отвечающую вопросу задачи.

  \item Проверьте знаки, единицы измерения и смысл округления.
\end{enumerate}

\section{Типичные ошибки}

\begin{enumerate}
  \item \warningtext{Начисление процентов на первоначальную сумму каждый год.}

  Неверно:
  \[
  S_3=S_0+3S_0p.
  \]

  При сложных процентах каждый год используется новая сумма:
  \[
  S_3=S_0(1+p)^3.
  \]

  \item \warningtext{Путаница между \(P\), \(p\) и \(r\).}

  При ставке \(12\%\):
  \[
  P=12,\qquad p=0{,}12,\qquad r=1{,}12.
  \]

  \item \warningtext{Начисление процентов на поздний взнос слишком рано.}

  Если деньги внесены после начисления процентов, в текущем году они процентов
  не приносят.

  \item \warningtext{Потеря знака при снятии денег.}

  Снятие \(x\RUB\) означает:
  \[
  x_k=-x.
  \]

  \item \warningtext{Неправильное раскрытие скобок после вычитания.}

  Если
  \[
  \Delta=A-(B-C+D),
  \]
  то
  \[
  \Delta=A-B+C-D.
  \]

  \item \warningtext{Обычное округление вместо смыслового.}

  Если минимальный срок равен \(8{,}13\) года, ответом является \(9\) полных
  лет, а не \(8\).

  \item \warningtext{Деление на неизвестную первоначальную сумму без проверки.}

  В задачах о вкладе обычно \(S_0>0\), поэтому сокращение на \(S_0\) допустимо.
  Это условие следует проговаривать.

  \item \warningtext{Применение формулы геометрической прогрессии при \(r=1\).}

  При нулевой ставке знаменатель \(r-1\) равен нулю. В этом случае используется
  формула
  \[
  S_n=S_0+nx.
  \]
\end{enumerate}

\section{Итоговый чек-лист}

Перед завершением решения проверьте:

\begin{itemize}
  \item записано ли \(p=P/100\);
  \item записано ли \(r=1+p\);
  \item установлен ли точный порядок начислений и операций;
  \item составлена ли таблица по всем периодам;
  \item каждое ли снятие имеет знак минус;
  \item переносится ли конечная сумма года в начало следующего года;
  \item начисляются ли проценты именно на актуальную сумму;
  \item правильно ли расставлены степени коэффициента \(r\);
  \item аккуратно ли раскрыты скобки;
  \item допустимы ли выполненные сокращения и деления;
  \item соответствует ли округление смыслу задачи;
  \item записан ли ответ с нужными единицами измерения.
\end{itemize}

% =================================================
% Тренировочный блок
% =================================================
\clearpage
\section*{Тренировочный блок}

Во всех задачах проценты начисляются один раз в конце каждого года.
Дополнительные операции выполняются сразу после начисления процентов.
Денежные ответы округляйте до копеек, если получается нецелое число.

\subsection*{Средний уровень}

\begin{enumerate}[label=\textbf{\arabic*.}]

  \item На вклад положили \(20\,000\RUB\) под \(8\%\) годовых.
  Дополнительных операций не было. Найдите сумму на счёте через три года.

  \item На вклад положили \(50\,000\RUB\) под \(10\%\) годовых.
  После каждого из двух начислений процентов вкладчик вносил ещё
  \(5000\RUB\). Найдите сумму на счёте через два года.

  \item На вклад положили \(100\,000\RUB\) под \(12\%\) годовых.
  После первого начисления процентов сняли \(10\,000\RUB\), после второго
  внесли \(5000\RUB\), после третьего дополнительных операций не выполняли.
  Найдите итоговую сумму.

\end{enumerate}

\subsection*{Уровень выше среднего}

\begin{enumerate}[label=\textbf{\arabic*.},resume]

  \item Первоначальная сумма вклада равна \(S\). Ставка составляет \(9\%\)
  годовых, дополнительных операций нет. Выразите сумму вклада через четыре
  года через \(S\).

  \item Вклад открыт под \(7\%\) годовых. Дополнительных операций нет.
  Через какое минимальное целое количество лет сумма станет не меньше
  удвоенной первоначальной суммы?

  \item Вклад открыт под \(13\%\) годовых. Дополнительных операций нет.
  Через какое минимальное целое количество лет сумма станет не меньше
  утроенной первоначальной суммы?

  \item Вклад открыт под \(12\%\) годовых на три года. После первого
  начисления процентов вкладчик снял \(3000\RUB\), после второго внёс
  \(3000\RUB\), после третьего операций не выполнял. На сколько рублей
  полученная сумма меньше той, которая была бы получена без промежуточных
  операций?

  \item На вклад положили \(80\,000\RUB\) под \(6\%\) годовых.
  После каждого из четырёх начислений процентов дополнительно вносили
  \(10\,000\RUB\). Найдите сумму на счёте после четвёртого взноса.

  \item Вклад открыт под \(10\%\) годовых на три года. После каждого
  начисления процентов вносили по \(15\,000\RUB\). Через три года на счёте
  оказалось \(200\,000\RUB\). Найдите первоначальную сумму вклада.

\end{enumerate}

\subsection*{Сложный уровень}

\begin{enumerate}[label=\textbf{\arabic*.},resume]

  \item Вклад в размере \(120\,000\RUB\) открыт под \(10\%\) годовых
  на четыре года. После первого начисления процентов вкладчик снял
  \(x\RUB\), после второго внёс \(2x\RUB\), после третьего снова снял
  \(x\RUB\), а после четвёртого операций не выполнял.

  В результате итоговая сумма оказалась на \(330\RUB\) меньше суммы,
  которая была бы получена без промежуточных операций. Найдите \(x\).

\end{enumerate}

% =================================================
% Ответы
% =================================================
\clearpage
\section*{Ответы к тренировочному блоку}

\begin{table}[ht]
\centering
\caption{Краткие ответы}
\begin{tabularx}{\linewidth}{@{}cX@{}}
\toprule
\textbf{№} & \textbf{Ответ} \\
\midrule

1 &
\(25\,194{,}24\RUB\). \\

2 &
\(71\,000\RUB\). \\

3 &
\(133\,548{,}80\RUB\). \\

4 &
\[
S_4=S\cdot1{,}09^4=1{,}41158161S.
\]
\\

5 &
\[
n=\left\lceil\frac{\lg2}{\lg1{,}07}\right\rceil
=\lceil10{,}244\ldots\rceil=11.
\]
Ответ: \(11\) лет. \\

6 &
\[
n=\left\lceil\frac{\lg3}{\lg1{,}13}\right\rceil
=\lceil8{,}988\ldots\rceil=9.
\]
Ответ: \(9\) лет. \\

7 &
\(403{,}20\RUB\). \\

8 &
\(144\,744{,}32\RUB\). \\

9 &
\(112\,960{,}18\RUB\). \\

10 &
\(x=30\,000\RUB\). \\

\bottomrule
\end{tabularx}
\end{table}

\vfill

\begin{center}
\color{darkaccent}
\textbf{Финансовая задача становится управляемой,\\
если строго соблюдать порядок операций и последовательно заполнять таблицу.}
\end{center}

\end{document}